\subsection{Generic recognition and the necessary profinite input}
\label{sec:archive-recognition}

The last sentence of scan~85 asks for recognition of the distinguished
geometric subtorsor. There are two different input diagrams here.
The isolated normalizer extension does not determine that subtorsor.
The full relative curve group together with the generic Galois marking
used in scans~83--84 does determine it, by the following argument.
The conclusion is functorial determination and an exact membership
criterion; it is not a finite-quotient decision algorithm.

\begin{theorem}[F149-GENERIC]\label{thm:archive-generic}
Retain the hypotheses and the normalizer identification of
Theorem~\ref{thm:archive-relative}. Suppose that $S$ is integral and
normal, with function field $K$ finitely generated over $\mathbb Q$.
Let $V$ be the normal line of the marked boundary section, let
$I\simeq\widehat{\mathbb Z}(1)$ be its geometric inertia, and put
\[
 G=\pi_1(S),\qquad G_K=\operatorname{Gal}(\overline K/K),\qquad
 N=\pi_1(V^\times).
\]
For an $I$-conjugacy class $s$ of sections of $N\to G$, the following
conditions are equivalent:
\begin{enumerate}[label=\textup{(\roman*)}]
\item $s$ is induced by a nowhere-zero global section of $V$;
\item the pullback $s_K$ to $G_K$ is induced by a point of
      $V_K^\times(K)$.
\end{enumerate}
The realizing global section, when it exists, is unique.
No finite-generation hypothesis on $\operatorname{Pic}(S)$ is needed
for this equivalence.
\end{theorem}

\begin{proof}
The map $G_K\to G$ is surjective: a connected finite \'etale cover
of a normal integral scheme has connected generic fibre. The
fundamental exact sequence over the field $K$, the assumed exact
sequence over $S$, and their common geometric fibre identify
\[
 N_K:=\pi_1(V_K^\times)\simeq N\mathbin{\times}_G G_K.
                                                               \tag{RG.1}
\]
Indeed the natural map of these two extensions induces the identity
on the kernel $I$ and on $G_K$, and is therefore an isomorphism.
The same argument gives
$\pi_1(U_K)\simeq\pi_1(U)\times_GG_K$ for the full curve groups.

A global nowhere-zero section clearly gives its generic section,
so only the converse needs proof. Suppose $s_K$ is induced by a
nonzero rational vector $v\in V_K(K)$. On an affine open
$W\subset S$ trivializing $V$, choose a basis $b$ and write
$v=xb$ with $x\in K^\times$. The basis gives a reference section
of the normalizer extension over $W$. Its difference from the
restriction of $s$ is a continuous $I$-valued cocycle. Reducing
modulo $n$ gives a class in $H^1(W,\mu_n)$, whose generic restriction
is the Kummer class of $x$, up to the fixed normal/conormal sign.
That sign has no effect on the argument below.

For each codimension-one point $w\in W$, the local ring
$A=\mathcal O_{S,w}$ is a discrete valuation ring. Restricting the
class to $\operatorname{Spec}A$ and using $\operatorname{Pic}(A)=0$
in the finite Kummer sequence shows that the class of $x$ in
$K^\times/K^{\times n}$ is represented by an element of $A^\times$.
Thus $x=a_n y_n^n$ for some $a_n\in A^\times$, $y_n\in K^\times$,
and
\[
 \operatorname{ord}_w(x)=n\operatorname{ord}_w(y_n)
 \quad\hbox{for every positive integer }n.
\]
An integer divisible by every positive integer is zero. Hence
$\operatorname{ord}_w(x)=0$ at every codimension-one point of $W$.
Normality now implies both $x$ and $x^{-1}$ belong to
$\mathcal O(W)$: a noetherian normal domain is the intersection
of its height-one valuation rings inside its fraction field.
Equivalently, apply the normal Hartogs extension property in
\cite[Theorem~5.10.5 and proof, pp.~115--116]{EGAIV2} to the
locally free rank-one sheaf and its dual. The vector $v$ therefore
extends as a nowhere-zero section over every such $W$, and these
extensions agree on overlaps because they agree in $K$.
They give $v\in V^\times(S)$.

The two resulting section classes agree generically. Their
difference is in $H^1(G,I)$, and inflation
\[
 H^1(G,I)\longrightarrow H^1(G_K,I)
\]
is injective. This follows directly from cocycles: the kernel of
$G_K\to G$ acts trivially on $I$, and a cocycle inflated from $G$
which becomes a coboundary on $G_K$ is already that same
coboundary on $G$. Thus the section induced by $v$ is exactly $s$.

Finally, if two global nonzero vectors give the same section,
their ratio $a\in\mathcal O(S)^\times\subset K^\times$ has zero
full Kummer class in $H^1(G_K,I)$. The full-completion injectivity
F149-FULLINJ applies to this finitely generated characteristic-zero
field, so $a=1$. This proves uniqueness. No Picard-group assertion
was used in the converse.
\end{proof}

\begin{proposition}[F149-MARKED-RECOGNITION]
\label{prop:archive-marked-recognition}
Suppose that, in addition to the full curve extension and marked
geometric cusp inertia, the profinite input retains the diagram
\[
 \pi_1(U)\longrightarrow G\longleftarrow G_K
     \longrightarrow G_k,                                  \tag{RG.2}
\]
where $k$ is an identified number field, $K/k$ is a finitely
generated regular extension, and the last arrow is the actual
restriction map. This is the generic Galois marking used in
cote~149, scans~83--84, together with its arithmetic base.
Then the subset of geometric tangent section classes in scan~85
is functorially determined by this marked profinite diagram.
The assertion uses the full profinite groups in \textup{(RG.2)},
not just an isolated prime quotient or the extension $1\to I\to N\to G\to1$.
\end{proposition}

\begin{proof}
We specify both the reconstruction and why it is independent of
choices. The known anabelian inputs are Mochizuki's Theorems~A
and~B: a hyperbolic curve over a sub-$p$-adic field is fully faithful
for open fundamental-group homomorphisms over the base Galois
group, and the analogous full profinite statement holds for
function fields of arbitrary dimension over such a base.
Their precise original locators are
\cite[Theorem~A, p.~3; Theorem~B, p.~4;
Theorem~16.5 and the following profinite argument, p.~86;
Theorem~17.1 and proof, p.~87]{MochizukiLocal1999}.
Only the isomorphism consequences are needed here.
Number fields and their finitely generated extensions are sub-$p$-adic:
a finite presentation embeds into a finitely generated extension
of a finite extension of $\mathbb Q_p$, after choosing an
embedding of the number field into an algebraic closure of
$\mathbb Q_p$.

First reconstruct $K$ from $G_K\to G_k$ within the category of
regular finitely generated extensions of the identified $k$.
Theorem~B says exactly that two realizations and any isomorphism
of their marked Galois groups differ by a unique field isomorphism.
For transcendence degree zero this is just $K=k$.
Form, entirely from \textup{(RG.2)}, the fibre product
\[
 \Pi_K=\pi_1(U)\mathbin{\times}_G G_K.
\]
By (RG.1) and its curve analogue, this is the full arithmetic
fundamental group of $U_K$ over $G_K$. Theorem~A reconstructs
$U_K$ up to its unique corresponding isomorphisms. The marked
inertia class determines its marked cusp in the smooth proper
compactification. Its normal line $V_K$, and the tangential
specialization map from its nonzero rational vectors to section
classes of $N_K$, are then reconstructed along with the curve.
Tangential specialization is functorial by the same relative
construction of Deligne used in Theorem~\ref{thm:archive-relative}.

Here is the explicit resulting membership criterion. Choose any
nonzero vector $b\in V_K(K)$ in the reconstructed one-dimensional
space, and let $s_b$ be its tangential section. For a proposed
global splitting $s$, form
\[
 \delta_b(s_K)=[s_K]-[s_b]\in H^1(G_K,I).
\]
Then
\[
 s\text{ is geometric over }S
 \quad\Longleftrightarrow\quad
 \delta_b(s_K)\in
 \kappa_K(K^\times),                                      \tag{RG.3}
\]
where $\kappa_K$ is the full Kummer map of the reconstructed field,
using the inertia orientation of the reconstructed marked curve.
For the forward implication apply Theorem~\ref{thm:archive-generic}.
For the reverse implication, a Kummer representative $a\in K^\times$
gives the rational vector $ab$ (or its inverse convention), and
Theorem~\ref{thm:archive-generic} extends it over $S$.
Replacing $b$ by $ub$, $u\in K^\times$, translates
$\delta_b$ by $\kappa_K(u)$, so does not change the condition.

It remains to check that the criterion has not concealed a choice
of geometric realization. Given two realizations of the marked
field group, Theorem~B supplies the corresponding field isomorphism.
After transport by it, two realizations of $\Pi_K\to G_K$ are
related by the curve isomorphism supplied by Theorem~A. This
isomorphism extends to the smooth proper compactifications, sends
the marked cusp to the marked cusp, and its derivative bijects
their nonzero rational normal vectors. Naturality of tangential
specialization identifies their sets of geometric section classes.
The group isomorphisms are defined up to geometric conjugacy.
Once the marked inertia representatives are aligned, the remaining
geometric conjugacy lies in $I$, because cusp inertia is
self-normalizing in the geometric curve group
\cite[Proposition~16, p.~8]{StixCuspidal2012}.
Thus the induced set of $I$-conjugacy classes is unambiguous.
Consequently every isomorphism of the full marked diagram
preserves (RG.3). This proves functorial determination.

This argument assumes that the marked diagram is known to arise
from the geometric situation stated in the proposition. It neither
recognizes arbitrary profinite groups as such diagrams nor supplies
an effective decision from finitely many finite quotients. It also
does not recover the additional generic marking from $G$ alone.
\end{proof}

\begin{proposition}[F149-BARE-EXTENSION]
\label{prop:archive-bare-extension}
The isolated extension $1\to I\to N\to G\to1$, even with its
specified inertia module and cyclotomic action, does not in general
intrinsically determine the geometric splitting subset. More
precisely, that subset need not be preserved by automorphisms of
the extension acting identically on $I$ and $G$.
\end{proposition}

\begin{proof}
Take $S=\operatorname{Spec}\mathbb Q$ and the trivial normal line,
for example at the cusp $0$ of the tripod. A reference tangent
vector identifies the normalizer with
\[
 N=I\rtimes G_{\mathbb Q},\qquad I=\widehat{\mathbb Z}(1).
\]
For every continuous cocycle $z:G_{\mathbb Q}\to I$, the map
\[
 (t,g)\longmapsto(t+z(g),g)                                 \tag{RG.4}
\]
is a continuous group automorphism, with inverse obtained by
replacing $z$ by $-z$. The cocycle identity verifies multiplication,
and the map fixes both the kernel and the quotient. On section
classes it acts by translation by $[z]\in H^1(G_{\mathbb Q},I)$.
Such translations are transitive on all section classes.

The geometric subset, relative to the chosen tangent vector, is
$\kappa_{\mathbb Q}(\mathbb Q^\times)$, which is proper.
For an explicit extra class take
$\alpha\in\widehat{\mathbb Z}$ whose $2$-adic component is
$-1/3\in\mathbb Z_2$ and whose other prime components are zero.
The compatible classes of $2^{\alpha\bmod n}$ in
$\mathbb Q^\times/\mathbb Q^{\times n}$ define a full Kummer class.
If it came from $a\in\mathbb Q^\times$, the ordinary integer
$\operatorname{ord}_2(a)$ would equal $\alpha$ in
$\widehat{\mathbb Z}$, hence would equal $-1/3$ in $\mathbb Z_2$,
which no integer does. Thus this class is not geometric.
The continuous-cohomology/Kummer comparison F149-KUM realizes it
by a continuous cocycle $z$, and (RG.4) sends the geometric
reference section to this nongeometric class.

This proves the assertion about the isolated extension. It does
not give an automorphism of the full curve diagram (RG.2):
Proposition~\ref{prop:archive-marked-recognition} shows precisely
why retaining that diagram removes this ambiguity.
\end{proof}

The scope of scan~85 can therefore be stated precisely. Geometric
splittings are recognizable from the full relative curve group
with its generic arithmetic marking, by (RG.3); generic geometricity
already forces global geometricity by normality and all-exponent
Kummer ramification. They cannot be selected from the isolated
normalizer extension by an automorphism-invariant rule. Recovering
the generic marking from a single unmarked base fundamental group
is not part of this positive theorem. The following subsection
refutes functorial recovery from the specified smaller full diagram
when its base group is allowed to move.
These statements are independent of step~17 and of the higher
$K$-group vanishing problem.
